\section{Industry-Standard Protocols}
\label{sec:industry-protocols}

This section applies the rubric of \S\ref{sec:rubric} to the
five inter-agent protocols that have either been adopted into
a formal standards body, ship with named production
deployments, or both: Google A2A, the Model Context Protocol,
IBM ACP, Cisco-led AGNTCY/SLIM, the Agent Network Protocol,
and Coral. The commerce-settlement industry-led
protocols (AP2, Visa TAP, Mastercard Agent Pay, ERC-8004, UCP) are
treated as a co-stack in \S\ref{sec:emerging:commerce} because
their design intent and rubric alignment differ from the
peer-messaging and tool-binding members of this group. Each
sub-section names the protocol's design intent, summarizes
the rubric scores, and identifies the open issues that
implementations and successor specifications must address.

\subsection{Google A2A}
\label{sec:industry:a2a}

A2A was introduced by Google in April 2025 as the canonical
peer-messaging protocol for inter-agent
coordination~\cite{spec_a2a}, and was donated to the Linux
Foundation under the AAIF umbrella in late 2025~\cite{lf_aaif}.
The protocol's defining structural choice is to treat the
\texttt{Task} as a first-class lifecycle-bearing object with
explicit submitted, working, input-required, auth-required,
completed, failed, and canceled states, and to pair the
task object
with a typed \texttt{Part} content model (TextPart,
FilePart, DataPart) that supports inline and by-reference
delivery. The discovery convention is a well-known endpoint
at \texttt{/.well-known/agent-card.json} returning a typed
AgentCard descriptor.

Rubric scores at the 2026-Q1 cutoff: Axis 1 = 1 (well-known
endpoint discovery), Axis 2 = 1 (OAuth-mediated bearer
identity), Axis 3 = 1 (advertised authentication schemes
and an in-task authorization flow, but no standardised
delegation credential), Axis 5 = 3 (full task object
with artifacts and sub-task delegation), Axis 6 = 3 (push
notifications combined with durable resumable tasks),
Axis 7 = 3 (typed parts with inline-or-reference model),
Axis 8 = 1 of 12 (transport encryption is the only
spec-level mitigation; bearer identity, push-notification
replay verification, and transport-layer DoS resistance are
partial and so do not count; injection, registry poisoning,
and delegation token theft are not addressed), Axis 9 = 1 (free-text Skill
descriptions over typed Skill objects), Axis 10
observability sub-score = 2 (SSE event stream supports
tracing; no standardized audit primitives).

The transport tuple is rich: HTTPS/JSON-RPC 2.0 as the
primary binding, optional gRPC and REST, JSON
serialization, SSE streaming, push notifications for async
delivery~\cite{spec_a2a}. The protocol composes naturally
with MCP for the tool-binding layer of each agent and with
AP2 or Visa TAP for the commerce-settlement layer
(\S\ref{sec:taxonomy:complementary}).

The specification does address authorization, and the
survey's Axis 3 score should not be read as saying
otherwise. \S7 states security requirements, the AgentCard
advertises authentication schemes, and \S7.6 defines
in-task authorization: an agent that needs further
authorization mid-task moves to the auth-required state and
delegates fulfilment of the request to its
client~\cite{spec_a2a}. What the specification leaves to
the surrounding identity architecture is authorization
policy (\S7.5 makes it implementation-specific) and the
choice of delegation mechanism. The score records that
absent credential format, not an absent security model.

The exclusion is deliberate and the specification says so
in normative language. \S7.6.4 states that A2A does not
define the scope, representation, validity, or revocation
semantics of the credential attached to the auth-required
state. It requires that the state transition alone must not
be treated as authorization for any operation, and that a
credential obtained through it must not be assumed to
authorize later messages on the same task. Meaning and scope
must come from the implementation, the credential issuer, or
an A2A extension~\cite{spec_a2a}. So A2A specifies
authorization \emph{knobs} rather than an authorization
model, and does it on purpose, so that many authorization
solutions can be plugged in above it. The survey reads that
as a scoping decision rather than an oversight. The Axis 3
score measures what a receiver can verify from the protocol
alone, which is why a deliberate exclusion and an accidental
omission land on the same level, and why
\S\ref{sec:security:authz} separates the two in prose.

The open issues against A2A at the cutoff are at the
identity and authorization layers. Eight named weaknesses
are documented in the security
literature~\cite{a2a_weaknesses}: AgentCard supply-chain
attacks, registry poisoning against the discovery convention,
delegation token theft, sensitive-data leakage through
free-text Skill descriptions, prompt-injection vectors via
DataPart payloads, replay attacks on push notifications,
side-channel privacy leakage from task lifecycle
transitions, and absent cryptographic delegation chains.
The A2ASecBench protocol-aware
evaluation~\cite{a2asecbench} catalogues several of these
empirically. Mitigations under discussion in the
Linux Foundation working group include the
ANP-style DID overlay for identity, the agentic-JWT
delegation pattern~\cite{agentic_jwt}, and the
authorization-propagation chain
proposed in~\cite{authz_propagation}; none has been
adopted into the spec as of 2026-Q1.

\subsection{Model Context Protocol (MCP)}
\label{sec:industry:mcp}

MCP is the canonical tool-binding
protocol~\cite{spec_mcp}: the asymmetric two-party
specification by which an agent (the client) invokes
external capabilities exposed as tools, resources, and
prompts. It was introduced by Anthropic in late 2024,
adopted by every major model vendor within months, and
governed under the Linux Foundation AAIF alongside
A2A~\cite{lf_aaif}.

Rubric scores: Axis 1 = 1 (server config-based discovery
with an emerging registry, mostly out-of-band), Axis 2 = 1
(OAuth 2.1 for HTTP transport, stdio relies on parent-process
trust), Axis 3 = 1 (OAuth scopes only, no agent-delegation
primitives), Axis 5 = 1 (request/response with sessions
but no protocol-level state machine; the emerging MCP Tasks
extension begins to add task semantics), Axis 6 = 2
(notifications plus emerging MCP Tasks; long-running
operations supported but not durably resumable across
connection loss), Axis 7 = 2 (typed content blocks for
text, image, and resource; structured JSON-Schema tool
input), Axis 8 = 1 of 12 (transport encryption is the
only spec-level mitigation; OAuth 2.1 identity is partial,
because the stdio transport falls back to process trust;
injection and tool poisoning remain
unaddressed at spec level~\cite{zhan2024injecagent,
protocol_exploits}), Axis 9 = 1 (JSON Schema for tool
input combined with free-text tool descriptions),
Axis 10 observability = 1 (logging hooks emerging; no
standardized audit primitives).

The transport tuple distinguishes MCP from A2A most
clearly: stdio, HTTP+SSE, and streamable-HTTP bindings; the
JSON-RPC 2.0 serialization shared with A2A; SSE or chunked
HTTP streaming; notifications and emerging MCP Tasks for
async delivery~\cite{spec_mcp}. The stdio binding is
particularly important in production: it allows a sandboxed
agent process to invoke a tool process within the same
trust boundary without the network attack surface of an
HTTP binding. The trade-off is that stdio inherits the
parent process's identity completely; the identity axis
score of 1 hides a sub-distinction between HTTP+OAuth and
stdio+process-trust that the Axis 4 tuple makes explicit.

The principal open issues at the cutoff are tool poisoning,
indirect prompt injection through tool results, and the
need for richer authorization primitives than OAuth
scopes can express. The InjecAgent
benchmark~\cite{zhan2024injecagent} establishes that
indirect injection is a measurable threat against deployed
MCP servers; the protocol does not currently specify
mitigations. The AAIF working group has begun discussing
an MCP-level authorization extension that would carry
agent-principal-delegation chains alongside OAuth scopes;
this would raise the Axis 3 score and align MCP with the
emerging chain-of-custody pattern that the
commerce-settlement layer requires.

\subsection{IBM ACP / BeeAI}
\label{sec:industry:acp}

IBM's Agent Communication Protocol (ACP), incubated under
the BeeAI project, occupies the tool-binding layer
alongside MCP but with structurally different
choices~\cite{spec_acp}. The transport is REST over HTTP,
in contrast to MCP's JSON-RPC 2.0. Serialization is JSON,
with streaming delivered as Server-Sent Events. The
positioning is stateless by design, in the sense that HTTP
is stateless: sessions are available when an agent needs
to carry state across runs.

ACP requires a status caveat that no other row in
Table~\ref{tab:rubric-matrix} needs. The specification is
\emph{complete but terminated}. The upstream repository was
archived in August 2025, and the project merged into A2A
under the Linux Foundation umbrella~\cite{spec_acp}. The
scores we report describe the specification as published,
not an effort still competing for adoption. We retain the
row because ACP is the survey's clearest evidence for
Finding~5 (\S\ref{sec:conclusion:governance}): it scores at
the top of the set on the task axis and still lost.

Rubric scores: Axis 1 = 1 (a well-known manifest at
\texttt{/.well-known/agent.yml} plus a live
\texttt{GET /agents} endpoint; registry-based discovery is
named but unspecified, and no descriptor is signed),
Axis 2 = 1 (the OpenAPI document declares no security
schemes at all; Basic, Bearer, and JWT appear only in prose,
enforced at a reverse proxy), Axis 3 = 1 (no authorization
primitive whatsoever), Axis 5 = 3 (a first-class
\texttt{Run} object with an eight-state lifecycle,
two-phase cancellation, and named artifacts), Axis 6 = 2
(asynchronous accept-and-poll, event backfill, and
distributed sessions portable across servers by URL
descriptor), Axis 7 = 2 (typed parts with inline-or-reference
duality and manifest-level content-type negotiation),
Axis 8 = 1 of 12 (transport encryption only, and that from
prose rather than specification text; the other eleven
threat classes are not addressed even partially, which is
unique in the set), Axis 9 = 1 (typed
manifest with free-text descriptions), Axis 10
observability = 2 (OpenTelemetry integration with default
exporters, plus trajectory and citation metadata carried in
the wire format).

Two features distinguish ACP, and neither is a content
format. The first is the run lifecycle. ACP defines an
\texttt{Await} mechanism by which an agent pauses execution
to request input from its client, then resumes on a
subsequent call. Combined with the eight-state machine and
named artifacts, this puts ACP at the same level as A2A on
the task axis and two levels above MCP. The second is
distributed sessions: session history and state are
referenced by URL on arbitrary resource servers, so a
session migrates between independent ACP servers by
forwarding only its descriptor, with no shared
infrastructure.

The gap between Axis 5 (level 3) and Axis 8 (1 of 12) is the
widest of any protocol we score, tied with A2A, and in ACP's
case it is deliberate rather
than accidental. ACP is HTTP-native by design, and its
answer to security is that the surrounding HTTP
infrastructure already provides one. That is a defensible
engineering position and a weak specification position. ACP
documents four multi-agent composition patterns, including
hierarchical delegation, while specifying no way for a
receiver to establish on whose authority a caller acts. It
is the sharpest instance in the survey of the
authorisation-propagation gap that Finding~3 names.

\subsection{AGNTCY / SLIM (Cisco-led)}
\label{sec:industry:agntcy}

AGNTCY is a Cisco-led collective effort hosted within the
Linux Foundation umbrella. It positions itself as a
``directory plus identity plus observability stack'' rather
than a single protocol, and that positioning has a
methodological consequence for this survey: AGNTCY is not
one specification but four. We score it across the Open
Agentic Schema Framework (OASF) for agent records, the
Agent Directory Service (ADS) for publication and
discovery, the Secure Low-latency Interactive Messaging
protocol (SLIM) for the messaging substrate, and a separate
identity specification~\cite{spec_agntcy}. ADS and SLIM are
each published as IETF Internet-Drafts. Both are individual
Independent Submissions with informational intended status,
so neither carries working-group adoption or IETF
endorsement.

OASF is the descriptor contribution, and its capability
model is stronger than free-text schemas. Records must
declare skills drawn from a published hierarchical
taxonomy, currently eighteen skill categories with nested
subcategories, alongside a parallel taxonomy of twenty
domain verticals. The framework is explicitly modelled on
OCSF, the Open Cybersecurity Schema Framework, and its
schema server is a derivative of the OCSF server. A
controlled vocabulary that constrains capability
descriptors to a shared term space is our Axis 9 level-2
anchor, so OASF scores 2. It does not reach level 3:
matching is taxonomic lookup, not semantic-similarity
search or logic-based reasoning producing verifiable
alignments.

ADS is the discovery contribution, and it is the only
industry-led effort in the survey to reach Axis 1 level 3.
Records are content-addressed by CID, so a descriptor's
identifier is a digest of its contents. Independent
directory nodes federate over a distributed hash table with
no central registry: publication announces
(CID, skills) pairs into the routing layer, and lookup
queries the skill taxonomy to retrieve matching CIDs and
their hosting nodes. Records are signed as
Cosign-compatible artifacts, with keyless OIDC, self-managed,
and KMS-backed signing flows all specified. Verifiable
names bind a signing key to a claimed domain by requiring a
JSON Web Key Set at that domain's well-known endpoint.
Discovery therefore returns cryptographically verifiable
descriptors resolvable independently of any single
authority, which is the level-3 descriptor.

SLIM is the messaging substrate, and its security
architecture is two-tier rather than transport-only. At the
network layer, HTTP/2 and HTTP/3 provide hop-by-hop TLS 1.3
encryption. HTTP/2 is the deliberate primary choice, for
compatibility with deployed load balancers, gateways, and
CDNs. Above it, Message Layer Security~\cite{rfc9420}
provides an application-layer encryption envelope
independent of the underlying TLS sessions. The consequence
is that routing nodes are zero-trust intermediaries: they
forward opaque encrypted blobs and cannot read message
content even after TLS terminates. SLIM names entities with
decentralized identifiers rather than transport
credentials. Node names are \texttt{did:key} values, client
locators compose an organization DID, namespace, service,
and client DID, and a channel is keyed to the DID of its
MLS group moderator, derived as a hash of that moderator's
public key.

Rubric scores: Axis 1 = 3 (content-addressed records over
a federated DHT with signed descriptors), Axis 2 = 2
(W3C Verifiable Credential agent badges verified through
resolver metadata, plus SPIFFE workload identity),
Axis 3 = 1 (MLS group membership plus short-lived OAuth
bearer tokens; no delegation chain), Axis 5 = 1 (sessions,
not tasks), Axis 6 = 3 (session state, group membership,
subscriptions, and pending queues survive disconnection and
node failure, with automatic reconnection and subscription
recovery), Axis 8 = 6 of 12, Axis 9 = 2 (the OASF
taxonomy), Axis 10 observability = 3, the highest of the
industry-led set, on the strength of a dedicated
observability sub-project shipping an OpenTelemetry SDK.

Two scores warrant explanation. Axis 2 sits at 2 rather
than 3 despite the use of Verifiable Credentials, because
the specification explicitly disclaims being bound by the
DID and VC standards, and its worked examples anchor agent
identifiers to enterprise identity providers whose keys
resolve from a tenant key-set endpoint. That is
issuer-anchored identity with signed credentials, not
self-sovereign control of keys. Axis 7 is 0, and the score
is an artefact of scope rather than a deficiency: a
substrate whose intermediaries cannot decrypt payloads has
no place for a content model, and MLS makes that opacity a
design guarantee. We report it as not applicable in
substance, consistent with design goal D1
(\S\ref{sec:rubric:goals}).

The AGNTCY scores make the composition argument concrete
rather than merely asserted. AGNTCY supplies discovery at
level 3, async durability at level 3, and the strongest
supply-chain posture in the survey, precisely where A2A
scores 1, 3, and partial. A2A supplies a task model and a
content model at level 3, precisely where AGNTCY scores 1
and 0. The two are complementary by construction: an
A2A-speaking agent, published as an OASF record in a
federated directory, addressed over SLIM. The open issue at
the cutoff is not whether they overlap but which stack an
adopter depends on at each layer, and the literature is not
yet mature enough to answer that.

\subsection{Agent Network Protocol (ANP)}
\label{sec:industry:anp}

ANP is the most identity-forward of the peer-messaging
protocols~\cite{spec_anp}. Its design center is the
proposition that agent identity should be decentralized,
cryptographically verifiable, and resolvable independently
of any single registry. The protocol builds on the W3C
DID and Verifiable Credentials specifications, and is the
only surveyed protocol that scores level 3 on both
Axis 1 (DID-resolver-based discovery) and Axis 2 (DIDs +
VCs for identity).

Rubric scores: Axis 1 = 3 (DID resolver plus Agent
Description Protocol), Axis 2 = 3 (DIDs and VCs as W3C
standards), Axis 3 = 2 (VC-scoped capabilities but
chain-of-custody patterns are not fully standardized),
Axis 5 = 1 (multi-turn but lightly specified at the
protocol layer), Axis 6 = 1 (polling primarily; no
push-notification standard), Axis 7 = 1 (text plus file
references), Axis 8 = 5 of 12 (identity, authentication,
integrity, audit, and non-repudiation are strong through
DIDs; operational threats are weaker, namely prompt
injection, registry poisoning, and capability escalation),
Axis 9 = 2 (JSON-LD ontology base, ADP enables semantic
discovery), Axis 10 observability = 1 (built-in
non-repudiation logs but no standardized audit dashboard).

The transport tuple is the lightest of the surveyed
peer-messaging protocols: HTTPS only, JSON-LD over JSON
serialization, no native streaming, polling for async
delivery. ANP's design accepts the latency cost of polling
in exchange for the decentralization properties that
push-notification-based architectures complicate. The
white paper~\cite{spec_anp} explicitly positions ANP for
peer-to-peer deployments where no central registry can be
assumed.

The Axis 5 and Axis 6 scores need a word on intent,
because the low numbers are a design choice rather than an
unfinished edge. The account below follows the
specification, and the ANP Community confirmed the reading
in correspondence with the authors. ANP does not treat a
task as a
protocol-level concept at all. Its first-class concepts are
Information and Interface: Information is data the agent
publishes, and Interface is an API it exposes, either
natural-language or structured. A task is then one form of
an Interface rather than an abstraction the protocol
defines. The stated reason is generality. The task concept
is useful in many settings but is not universal enough to
carry every agent interaction, so ANP keeps its primary
abstractions broader. Asynchrony follows the same rule: an
Interface defines its own asynchronous model where it needs
one, and the protocol layer defines no separate primitive.
This is the same deliberate-exclusion pattern the survey
records for A2A on Axis 3
(\S\ref{sec:security:authz}). The axis measures what the
protocol itself specifies, so a considered omission and an
unfinished one land on the same level, and the survey
separates them in prose rather than in the number.

The open issues at the cutoff are therefore operational
maturity and the consequence of that scoping choice rather
than the choice itself. ANP's strengths are at the identity
and authentication layers; the dimensions it leaves to the
Interface are precisely those on which A2A scores highest
(tasks, async, and multi-modal handling). The natural
composition pattern (ANP for identity, A2A for task
semantics) would require bridging machinery the cutoff
specifications do not yet provide.\footnote{Post-cutoff, and
recorded here because it bears on Axis 6 and Axis 9 rather
than to revise them. ANP published a 1.1 release line in
June 2026, after the 2026-Q1 cutoff, adding a released
end-to-end instant messaging suite (ANP-09) for direct and
group agent-to-agent messages with attachments and optional
end-to-end encryption~\cite{spec_anp_v11}. Its premise is
that identity plus messaging carries a large share of
collaboration without modelling each interaction as a
structured task, which is the Information-and-Interface
position above followed to its conclusion. The same release
narrows the JSON-LD commitment: the Agent Description
document (ANP-07) moved to plain JSON for readability,
while the Discovery specification (ANP-08) keeps JSON-LD
and its \texttt{@context} mechanism. The community
describes the underlying Linked Data philosophy as
retained. Section~\ref{sec:semantic:ontology} records what
this would do to the Axis 9 reading.}

\subsection{Coral Protocol}
\label{sec:industry:coral}

Coral~\cite{spec_coral} occupies the peer-messaging layer
with two distinguishing choices: persistent multi-participant
threads as the conversational primitive, and an explicit
integration with MCP for the tool-binding layer. Coral
positions itself as ``the Internet of Agents'' substrate
and supports both off-chain coordination (over standard
HTTP and SSE) and on-chain settlement (via Solana wallets
and signed team contracts) where the deployment requires
it.

Rubric scores: Axis 1 = 2 (\texttt{list\_agents} MCP tool
resolves from a server-mediated Coral Server;
blockchain-anchored but discovery is server-mediated
rather than DID-resolver), Axis 2 = 3 (each agent has a
DID and cryptographic credentials with Solana wallets as
authenticator), Axis 3 = 2 (Team Contracts as signed
multi-agent agreements, optionally on-chain; not RFC 8693,
not a fully specified delegation chain), Axis 5 = 2
(stateful persistent multi-participant threads; sub-task
delegation deferred to a future Task Management service),
Axis 6 = 2 (\texttt{wait\_for\_mentions} is an explicit
push primitive; no durable resumable spec), Axis 7 = 0
(natural-language message content with mention metadata
only; no typed parts or content negotiation), Axis 8 = 3
of 12 (DIDs and wallets for auth, on-chain immutable
audit, and thread-scoped privilege confusion are
addressed; team-contract capability escalation,
delegation token theft, and signed-transaction integrity
are partial, the last because signatures cover the payment
path rather than the message path; replay, supply-chain,
side-channel, DoS, prompt injection, and registry
poisoning not addressed), Axis 9 = 1 (Coraliser adapters
enforce uniform schemas but no formal capability
descriptor or ontology), Axis 10 observability = 2
(blockchain audit trail is mandatory for payments and
team contracts; no separate trace or metric spec).

The transport tuple is HTTP, WebSocket, and SSE
bindings; JSON via MCP serialization; SSE streaming; push
via \texttt{wait\_for\_mentions}~\cite{spec_coral}. The
choice to deliver tool semantics through MCP rather than
re-specifying them is structurally important: Coral
explicitly accepts the layer composition argument of
\S\ref{sec:taxonomy:complementary} and builds on top of
MCP rather than alongside it.

The open issues at the cutoff are content-model
poverty (Axis 7 = 0 is the lowest score in the surveyed
set), the absent task lifecycle (Axis 5 stops at level 2),
and the dependency on blockchain primitives for the
strongest security claims. Deployments that do not want
to commit to Solana settlement infrastructure inherit a
weaker version of Coral's threat model.
